\section{Trust and reassurance: why this technology can be trusted}
\label{sec:trust}

Medical AI bills face honest public skepticism, so the seventh appeal is
reassurance. Trust is not asserted; it is built from named supports: physician
endorsement, academic review, mandatory safety testing, transparency provisions,
pilot results, and an audit trail that anyone can inspect. The objective is a
committee that can tell constituents, with evidence, that this technology can be
trusted.

[DRAFTING INSTRUCTIONS: in the full-narrative, build the trust scaffold from
\texttt{cancer-automated/papers/VVUQ-02} (the ten-gate suite, 172 of 172 tests,
each gate bound to a published consensus standard) and
\texttt{review/references/narrative\_refs.bib} entry \texttt{asme2018vv40} (ASME
V\&V 40-2018). Reproduce Mermaid figures
\texttt{review/mermaid/diagrams/10-trust-scaffolding.md},
\texttt{07-ten-vvuq-gates.md}, and \texttt{17-system-context.md} as colored TikZ.
Add a body-width table binding each gate to its standard and its pass threshold,
modeled on the VVUQ table convention.]

\begin{narrativefig}
\node[nfone] (a) at (0,0) {Standards};
\node[nfdec] (b) at (3.6,0) {Ten gates};
\node[nftwo] (c) at (7.2,0) {Audit trail};
\node[nfhope] (d) at (10.4,0) {Trust};
\draw[nfedge] (a)--(b); \draw[nfedge] (b)--(c); \draw[nfedge] (c)--(d);
\end{narrativefig}
\vspace{-0.2cm}
\figcaption{Figure 8. Trust built from standards and an audit trail (placeholder;
expanded from Mermaid 10 in the full-narrative).}

Reassurance works when it is specific. A named standard and a published test
result persuade where a promise of safety cannot.
